% !Mode:: "TeX:UTF-8"
% !TEX program = xelatex
\documentclass[a4paper]{article}
\usepackage{amsmath}
\usepackage{amssymb}
\usepackage{ctex}
%\usepackage{braket}
%\usepackage[european]{circuitikz}
\usepackage{multirow}
\usepackage{float}
\usepackage{graphicx}
\usepackage{geometry}
\geometry{left=2.5cm,right=2.5cm,bottom=2.5cm,top=2.5cm}
\title{近代物理实验报告4.4：高临界温度超导体临界温度的电阻测量法}
\author{xy\quad 学号\quad 匡亚明学院}
\date{2019年2月29日}
\begin{document}
\maketitle
\bibliographystyle{unsrt}
%--------main-body------------

\section{引言}
早在1991年荷兰物理学家卡麦林$\cdot$翁纳斯(Kamerlingh-Onnes)发现，将水银冷却到稍低于4.2K时，其电阻急剧地下降到零。他认为，这种电阻突然消失的现象，是由于物质转变到了一种新的状态，并将此以零电阻为特征的金属态，命名为超导态。1993年迈斯纳(Meissner)和奥森菲尔德(Ochsenfeld)发现超导电性的另一特性：超导态时磁通密度为零或叫完全抗磁性，即meissner效应。电阻为零及完全抗磁性是超导电性的两个最基本的特性。超导体从具有一定电阻的正常态，转变为电阻为零的超导态时，所处的温度叫做临界温度，常用$T_c$表示。直至1986年以前，人们经过70多年的努力才获得了最高临界温度为23K的Nb$_3$Ge超导材料。1986年4月，Bednorz和Muller 创造性地提出了在Ba-La-Cu-O系化合物中存在高$T_c$超导的可能性。1987年初，中国科学院物理研究所赵忠贤等在这类氧化物中发现了$T_c$=48K的超导电性。同年2月份，美籍华裔科学家朱经武在Y-Ba-Cu-O系中发现了$T_c$=90K的超导电性。这些发现使人们梦寐以求的高温超导体变成了现实的材料，可以说这是科学史上又一次重大的突破。其后，在1988年1月，日本科学家Hirashi Maeda报导研制出临界温度为106K的Bi-Ba-Ca-Cu-O系新型高温超导体。同年2月，美国阿肯萨斯大学的Allen Hermann和Z.Z.Sheng等发现了临界温度为106K的Tl-Ba-Ca-Cu-O系超导体。一个月后，IBM的Almanden又将这种体系超导体的临界温度提高到了125K。1989年5月，中国科技大学的刘宏宝等通过用Pb和Sb对Bi的部分取代，使Bi-Sr-Ca-Cu-O系超导材料的临界温度提高到了130K。这是迄今所报导的最高的临界温度。氧化物超导材料成材困难、韧性差、临界电流密度低，临界磁场不高等缺点局限了高温超导材料的应用，但已有包银铋锶钙铜氧组成线材获得越来越多的应用。

\section{实验目的}
\begin{enumerate}
\item 分别利用动态法和稳态法测量高临界温度氧化物超导材料的电阻率随温度的变化关系。
\item 通过实验掌握利用液氮容器内的低温空间改变氧化物超导材料温度、测温及控温的原理和方法。
\item 学习利用四端子法测量超导材料电阻和热电势的消除等基本实验方法以及实验结果的分析与处理。
\end{enumerate}

\section{实验仪器}
低温恒温器、测量仪器。

\section{实验原理}
\subsection{临界温度$T_c$的定义及其规定}
超导体具有零电阻效应，通常把外部条件(磁场、电流、应力等)维持在足够低值时电阻突然变为零的温度称为超导临界温度。实验表明，超导材料发生正常$\to$超导转变时，电阻的变化是在一定的温度间隔中发生，而不是突然变为零的，如图(\ref{fig1})所示。起始温度$T_s$(Onset Point)为R-T曲线开始偏离线性所对应的温度；中点温度$T_m$(Mid Point)为电阻下降至起始温度电阻$R_s$的一半时的温度；零电阻温度Ｔ为电阻降至零时的温度。而转变宽度$\Delta T$定义为$R_s$下降到90\%及10\%所对应的温度间隔。高$T_c$材料发现之前，对于金属、合金及化合物等超导体，长期以来在测试工作中，一般将中点温度定义为$T_c$，即$T_c$＝$T_m$。对于高$T_c$氧化物超导体，由于其转变宽度$\Delta T$较宽，有些新试制的样品$\Delta T$可达十几K，再沿用传统规定容易引起混乱。因此，为了说明样品的性能，目前发表的文章中一般均给出零电阻温度T(R=0)的数值，有时甚至同时给出上述的起始温度、中点温度及零电阻温度。而所谓零电阻在测量中总是与测量仪表的精度、样品的几何形状及尺寸、电极间的距离以及流过样品的电流大小等因素有关，因而零电阻温度也与上述诸因素有关、这是测量时应予注意的。
\begin{figure}[!h]
\centering
\includegraphics[width=0.6\textwidth]{fig/fig1.pdf}
\caption{超导材料地电阻温度曲线}\label{fig1}
\end{figure}

\subsection{样品电极的制作}
目前所研制的高$T_c$氧化物超导材料多为质地松脆的陶瓷材料，即使是精心制作的电极，电极与材料间的接触电阻也常达零点几欧姆，这与零电阻的测量要求显然是不符合的。为消除接触电阻对测量的影响，常采用图(\ref{fig2})所示的四端子法。两根电流引线与直流恒流电源相连，两根电压引线连至数字电压表或经数据放大器放大后接至X-Y记录仪，用来检测样品的电压。按此接法，电流引线电阻及电极1、4与样品的接触电阻与2、3端的电压测量无关。2、3两电极与样品间存在接触电阻，通向电压表的引线也存在电阻，但是由于电压测量回路的高输入阻抗特性，吸收电流极小，因此能避免引线和接触电阻给测量带来的影响。按此法测得电极2、3端的电压除以流过样品的电流，即为样品电极2、3端间的电阻。
\begin{figure}[!h]
\centering
\includegraphics[width=0.6\textwidth]{fig/fig2.pdf}
\caption{四端子接线}\label{fig2}
\end{figure}

\subsection{温度控制及测量}
临界温度$T_c$的测量工作取决于合理的温度控制及正确的温度测量。目前高$T_c$氧化物超导材料的临界温度大多在60K以上，因而冷源多用液氮。纯净液氮在一个大气压下的沸点为77.348K，三相点为63.148K，但在实际使用中由于液氮的不纯，沸点稍高而三相点稍低(严格地说，不纯净的液氮不存在三相点)。对三相点和沸点之间的温度，只要把样品直接浸入液氮，并对密封的液氮容器抽气降温，一定的蒸汽压就对应于一定的温度。在77K以上直至300K，常采用如下两种基本方法。
\subsubsection{普通恒温器控温法}
低温恒温器通常是指这样的实验装置。它利用低温流体或其他方法，使样品处在恒定的或按所需方式变化的低温温度下，并能对样品进行一种或多种物理量的测量。这里所称的普通恒温器控温法，指的是利用一般绝热的恒温器内的锰铜线或镍铬线等绕制的电加热器的加热功率来平衡液池冷量，从而控制恒温器的温度稳定在某个所需的中间温度上。改变加热功率，可使平衡温度升高或降低。由于样品及温度计都安置在恒温器内并保持良好的热接触，因而样品的温度可以严格控制并被测量。这样控温方式的优点是控温精度较高，温度的均匀性较好，温度的稳定时间长。用于电阻法测量时，可以同时测量多个样品。由于这种控温法是点控制的，因此普通恒温器控温法应用于测量时又称定点测量法。

由于恒温器与低温液体直接接触，有较多的冷量需较大的功率来平衡，同时各点的温度难以平衡，为了解决这个问题，只要减小冷量就可以了，一个简便的方法是在恒温器外包一层保温材料来制成绝热恒温器就可简单地控温。
\subsubsection{温度梯度法}
这是指利用贮存液氮的杜瓦容器内液面以上空间存在的温度梯度来自然获取中间温度的一种简便易行的控温方法。样品在液面以上不同位置获得不同温度。为正确反映样品的温度，通常要设计一个紫铜均温块，将温度计和样品与紫铜均温块进行良好的热接触。紫铜块连结至一根不锈钢管，借助于不锈钢管进行提拉以改变温度。

本实验的恒温器设计综合上述两种基本方法，既能进行动态测量，也能进行定点的稳态测量，以便进行两种测量方法和测量结果的比较。
\subsection{热电势及热电势的消除}
用四端子法测量样品在低温下的电阻时常会发现，即使没有电流流过样品，电压端也常能测量到几微伏至几十微伏的电压降。而对于高$T_c$超导样品，能检测到的电阻常在$\text{10}^{-5}\sim\text{10}^{-1}\Omega$之间，测量电流通常取100$\mu$A至100mA左右，取更大的电流将对测量结果有影响。据此换算，由于电流流过样品而在电压引线端产生的电压降只在$\text{10}^{-2}\mu\text{V}\sim\text{10}^3\mu\text{V}$之间，因而热电势对测量的影响很大，若不采取有效的测量方法予以消除，有时会将良好的超导样品误作非超导材料，造成错误的判断。

测量中出现的热电势主要来源于样品上的温度梯度。为什么放在恒温器上的样品会出现温度的不均匀分布呢？这取决于样品与均温块热接触的状况。若样品简单地压在均温块上，样品与均温块之间的接触热阻较大。同时样品本身有一定的热阻也有一定的热容。当均温块温度变化时，样品温度的弛豫时间与上述热阻及热容有关，热阻及热容的乘积越大，弛豫时间越长。特别在动态测量情形，样品各处的温度弛豫造成的温度分布不均匀不能忽略。即使在稳态的情形，若样品与均温块之间只是局部热接触(如不平坦的样品面与平坦的均温块接触)，由引线的漏热等因素将造成样品内形成一定的温度梯度。样品上的温差$\Delta T$会引起载流子的扩散，产生热电势Ｅ
\begin{equation}
E = S\Delta T\label{eq1}
\end{equation}
S是样品的微分热电势，其单位是$\mu\text{VK}^{-1}$。

对高$T_c$超导样品热电势的讨论比较复杂，它与载流子的性质以及电导率在费密面上的分布有关，利用热电势的测量可以获知载流子性质的信息。对于同时存在两种载流子的情况，它们对热电势的贡献要乘一权重，满足所谓Nordheim-Gorter法则。
\begin{equation}
S = \frac{\sigma_A}{\sigma}S_A + \frac{\sigma_B}{\sigma}S_B\label{eq2}
\end{equation}
式中$S_A$、$S_B$是A、B两种载流子本身的热电势，$\sigma_A$、$\sigma_B$分别为A、B两种载流子相应的电导率。$\sigma = \sigma_A+\sigma_B$。材料处在超导态时，S=0。

为消除热电势对测量电阻率的影响，通常采取下列措施：
\subsubsection{对于动态测量}
应将样品制得薄而平坦。样品的电极引线尽量采用直径较细的导线，例如直径小于0.1mm的铜线。电极引线与均温块之间要建立较好的热接触，以避免外界热量经电极引线流向样品。同时样品与均温块之间用导热良好的导电银浆粘接，以减少热弛豫带来的误差。另一方面，温度计的响应时间要尽可能小，与均温块的热接触要良好，测量中温度变化应该相对地较缓慢。对于动态测量中电阻不能下降到零的样品，不能轻易得出该样品不超导的结论，而应该在液氮温度附近，通过后面所述的电流换向法或通断法检查。
\subsubsection{对于稳态测量}
当恒温器上的温度计达到平衡值时，应观察样品两侧电压电极间的电压降及叠加的热电势值是否趋向稳定，稳定后可以采用如下方法。
\begin{enumerate}
\item 电流换向法：\\
将恒流电源的电流I反向，分别得到电压测量值$U_A$、$U_B$，则超导材料测电压电极间的电阻为
\begin{equation}
R = \frac{|U_A-U_B|}{2I}\label{eq3}
\end{equation}
\item 电流通断法：\\
切断恒流电源的电流，此时测电压电极间量到的电压即是样品及引线的积分热电势，通电流后得到新的测量值，减去热电势即是真正的电压降。若通断电流时测量值无变化，表明样品已经进入超导态。
\end{enumerate}

\section{实验内容}
\begin{enumerate}
\item 利用稳态法，在样品的零电阻温度与０℃之间测出样品的Ｒ -Ｔ分布。
\item 对实验数据进行处理、分析。
\item 对实验结果进行讨论。
\end{enumerate}

\section{注意事项}
\begin{enumerate}
\item 动态法测量时，热弛豫对测量的影响很大。它对热电势的影响随升降温速度变化以及相变点的出现可能产生不同程度的变化。应善于利用实验条件、观察热电势的影响。
\item 动态法测量中样品温度与温度计温度难以一致，应观察不同的升降温速度对这种不一致的影响。
\item 进行稳态法测量时可以选择样品在液面以上的合适高度作为温度的粗调，而以电脑给定值作为温度的细调。
\end{enumerate}

\section{实验数据}

\section{误差分析}

\section{思考题}
\subsection{超导样品的电极为什么一定要制作成如图(\ref{fig2})所示的四端子接法？假定每根引线的电阻为0.1$\Omega$，电极与样品间的接触电阻为0.2$\Omega$，数字电压表内阻为10M $\Omega$，试用等效电路分析当样品进入超导态时，直接用万用表测量与采用图(\ref{fig2})接法测量有何不同？}
\subsection{设想一下，本实验适宜先做动态法测量还是稳态法测量？为什么？}
\subsection{本实验的动态法升降温过程获得的R-T曲线有哪些具体差异？为什么会出现这些差异？}
\subsection{给出实验所用样品的超导起始温度、中间温度和零电阻温度，分析实验的精度。}
\subsection{实验中如何保证温度计测量的温度与样品的温度是相同的？}

\nocite{jiaocai}
\bibliography{ref}
\end{document}